\documentclass{article}
\usepackage[UTF8]{ctex}
\usepackage{graphicx}
\usepackage{fontspec}
\usepackage{xcolor}
\usepackage{amsmath}
\usepackage{amssymb}
\usepackage{bm}
\usepackage{amsfonts} 
\usepackage{tabularx}
\usepackage{booktabs}
\begin{document}
\title{Lecture Notes 7}
\author{Minfei Chen}
\date{\today}
\maketitle
\section{TD Algorithm}
TD 学习算法的更新规则如下：
\begin{align*}
    v_{t+1}(s_t) &= v_t(s_t) - \alpha_t(s_t) \left[ v_t(s_t) - [r_{t+1} + \gamma v_t(s_{t+1})] \right] \\
    v_{t+1}(s) &= v_t(s), \quad \forall s \neq s_t
\end{align*}
其中 $t = 0, 1, 2, \dots$。这里的 $v_t(s_t)$ 是对真实价值 $v_{\pi}(s_t)$ 的估计值；$\alpha_t(s_t)$ 是在时间 $t$ 时状态 $s_t$ 的学习率。

\subsubsection*{TD 更新规则分析}

TD 更新规则：
\[
\underbrace{v_{t+1}(s_t)}_{\text{新估计值}} = \underbrace{v_t(s_t)}_{\text{当前估计值}} - \alpha_t(s_t) \underbrace{\left[ v_t(s_t) - \overbrace{[r_{t+1} + \gamma v_t(s_{t+1})]}^{\text{TD 目标 } \bar{v}_t} \right]}_{\text{TD 误差}}
\]
 \textbf{TD 目标 (TD target)} 定义为：
\[
\bar{v}_t \doteq r_{t+1} + \gamma v_t(s_{t+1})
\]
“目标”的含义是，当前的估计值 $v_t(s_t)$ 会被“拉向”TD 目标 $\bar{v}_t$。

将 TD 目标代入更新规则：
\begin{align*}
    v_{t+1}(s_t) &= v_t(s_t) - \alpha_t(s_t) [v_t(s_t) - \bar{v}_t] \\
    % 两边同时减去 TD 目标
    \implies v_{t+1}(s_t) - \bar{v}_t &= v_t(s_t) - \bar{v}_t - \alpha_t(s_t) [v_t(s_t) - \bar{v}_t] \\
    % 合并同类项
    \implies v_{t+1}(s_t) - \bar{v}_t &= [1 - \alpha_t(s_t)][v_t(s_t) - \bar{v}_t] \\
    % 两边取绝对值
    \implies |v_{t+1}(s_t) - \bar{v}_t| &= |1 - \alpha_t(s_t)| |v_t(s_t) - \bar{v}_t|
\end{align*}
由于学习率 $\alpha_t(s_t)$ 是一个小的正数，我们有：
\[
0 < 1 - \alpha_t(s_t) < 1
\]
因此，我们可以得到：
\[
|v_{t+1}(s_t) - \bar{v}_t| \le |v_t(s_t) - \bar{v}_t|
\]
这证明了新的估计值 $v_{t+1}(s_t)$ 确实被拉向了 TD 目标 $\bar{v}_t$！

\textbf{TD 误差} $\delta_t$ 定义为：
\[
\delta_t = v(s_t) - [r_{t+1} + \gamma v(s_{t+1})]
\]

\begin{itemize}
    \item 它是两个连续时间步之间估计值的差异。即在 $t$ 时刻的估计 $v(s_t)$ 与在 $t+1$ 时刻得到的一个更可靠的目标值 $[r_{t+1} + \gamma v(s_{t+1})]$ 之间的差异。
    
    \item 它反映了当前估计值 $v_t$ 与真实价值 $v_{\pi}$ 之间的**“不足”或“缺陷” (deficiency)**。为了理解这一点，我们定义一个理论上的“贝尔曼误差”：
    \[
    \delta_{\pi,t} \doteq v_{\pi}(s_t) - [r_{t+1} + \gamma v_{\pi}(s_{t+1})]
    \]
    注意到，这个贝尔曼误差的期望值为零：
    \begin{align*}
        \mathbb{E}[\delta_{\pi,t} | S_t = s_t] &= v_{\pi}(s_t) - \mathbb{E}[R_{t+1} + \gamma v_{\pi}(S_{t+1}) | S_t = s_t] \\
        &= 0.
    \end{align*}
    
    \item 因此，如果估计值是完美的，即 $v_t = v_{\pi}$，那么我们计算出的 TD 误差 $\delta_t$ 的期望值也应该为零。
\end{itemize}

\begin{table}[htbp] % h: here, t: top, b: bottom, p: page of floats
    \centering
    \caption{TD 学习与 MC 学习的对比}
    \label{tab:td_vs_mc_full}
    \begin{tabularx}{\textwidth}{XX} % 使用两个X列，使它们宽度相等且自动换行
        \toprule
        \textbf{TD/Sarsa 学习} & \textbf{蒙特卡洛 (MC) 学习} \\
        \midrule
        \textbf{在线学习 (Online):} TD 学习是在线的。它可以在收到一个奖励后，立即更新状态或动作的价值。 & \textbf{离线学习 (Offline):} MC 学习是离线的。它必须等到一条经验片段 (episode) 被完整收集后，才能进行更新。 \\
        
        \addlinespace % 增加行间距
        
        \textbf{持续性任务 (Continuing tasks):} 由于 TD 学习是在线的，它可以处理有终点的分幕式任务，也可以处理没有终点的持续性任务。 & \textbf{分幕式任务 (Episodic tasks):} 由于 MC 学习是离线的，它只能处理那些有终止状态的分幕式任务。 \\
        
        \addlinespace
        
        \textbf{自举 (Bootstrapping):} TD 学习是自举的，因为其价值的更新依赖于之前的价值估计值。因此，它需要一个初始猜测值来启动。 & \textbf{非自举 (Non-bootstrapping):} MC 学习不是自举的，因为它可以直接估计状态/动作的价值而无需任何初始猜测值（它直接计算真实的回报）。 \\
        
        \addlinespace
        
        \textbf{低估计方差 (Low estimation variance):} TD 的方差比 MC 低，因为它只涉及较少的随机变量。例如，Sarsa 的更新只依赖于下一步的随机结果 $R_{t+1}, S_{t+1}, A_{t+1}$。 & \textbf{高估计方差 (High estimation variance):} 为了估计 $q_\pi(s_t, a_t)$，我们需要对一个很长的随机序列 $R_{t+1} + \gamma R_{t+2} + \dots$ 进行采样。路径中的任何一步随机性都会累积，导致回报值的波动很大。 \\
        \bottomrule
    \end{tabularx}
\end{table}
\section{Sarsa Algorithm}
\[
q_{t+1}(s, a) =
\begin{cases}
    q_t(s_t, a_t) - \alpha_t(s_t, a_t) \left[ q_t(s_t, a_t) - [r_{t+1} + \gamma q_t(s_{t+1}, a_{t+1})] \right], & \text{if } (s,a) = (s_t, a_t) \\
    q_t(s, a), & \forall (s, a) \neq (s_t, a_t)
\end{cases}
\]
其中 $t=0, 1, 2, \dots$
\begin{itemize}
    \item $q_t(s_t, a_t)$ 是对真实动作价值 $q_{\pi}(s_t, a_t)$ 的一个估计；
    \item $\alpha_t(s_t, a_t)$ 是与状态-动作对 $(s_t, a_t)$ 相关的学习率。
\end{itemize}

\subsubsection*{第一步：更新 Q 值 (Update q-value)}
根据采集到的经验五元组 $(s_t, a_t, r_{t+1}, s_{t+1}, a_{t+1})$，进行一次 Sarsa 更新：
\[
q_{t+1}(s_t, a_t) = q_t(s_t, a_t) - \alpha_t(s_t, a_t) \left[ q_t(s_t, a_t) - [r_{t+1} + \gamma q_t(s_{t+1}, a_{t+1})] \right]
\]

\subsubsection*{第二步：更新策略 (Update policy)}
基于刚刚更新得到的 Q 值，改进状态 $s_t$ 下的策略 $\pi$ （使用$\epsilon$-greedy）：
\[
\pi_{t+1}(a|s_t) =
\begin{cases}
    1 - \varepsilon + \frac{\varepsilon}{|A(s_t)|} & \text{如果 } a = \underset{a'}{\arg\max} \, q_{t+1}(s_t, a') \\
    \frac{\varepsilon}{|A(s_t)|} & \text{对于其他动作}
\end{cases}
\]
\section{expected Sarsa and n-step Sarsa}
\subsection{expected Sarsa}
expected Sarsa对 TD target 进行了改进：
\begin{align*}
    q_{t+1}(s_t, a_t) &= q_t(s_t, a_t) - \alpha_t(s_t, a_t) \left[ q_t(s_t, a_t) - (r_{t+1} + \gamma \mathbb{E}[q_t(s_{t+1}, A)]) \right], \\
    q_{t+1}(s, a) &= q_t(s, a), \quad \forall (s, a) \neq (s_t, a_t),
\end{align*}
其中，
\[
\mathbb{E}[q_t(s_{t+1}, A)] = \sum_{a} \pi_t(a|s_{t+1}) q_t(s_{t+1}, a) \doteq v_t(s_{t+1})
\]
这个期望值是 $q_t(s_{t+1}, a)$ 在当前策略 $\pi_t$ 下的期望。

\textcolor{blue}{\kaishu*Expected Sarsa通过计算某种策略下行动A的期望，从而减少了随机变量的个数，降低了算法的随机性，提高了收敛的速度}
\subsection{n-step Sarsa : unifying Sarsa and MC}
不同算法实际上是在使用不同“步数”的回报 (Return) 作为更新目标：
\begin{align*}
    \text{Sarsa (单步)} \quad & G_t^{(1)} = R_{t+1} + \gamma q_{\pi}(S_{t+1}, A_{t+1}), \\
    \quad & G_t^{(2)} = R_{t+1} + \gamma R_{t+2} + \gamma^2 q_{\pi}(S_{t+2}, A_{t+2}), \\
    &\vdots \\
    \text{n-步 Sarsa} \quad & G_t^{(n)} = R_{t+1} + \gamma R_{t+2} + \dots + \gamma^n q_{\pi}(S_{t+n}, A_{t+n}), \\
    &\vdots \\
    \text{蒙特卡洛 (MC)} \quad & G_t^{(\infty)} = R_{t+1} + \gamma R_{t+2} + \gamma^2 R_{t+3} + \dots
\end{align*}
\textcolor{blue}{\kaishu*n-step Sarsa 是一种特殊的online的方法}
\section{Q-learning}
Q-learning 算法的更新规则如下：
\[
q_{t+1}(s, a) =
\begin{cases}
    q_t(s_t, a_t) - \alpha_t(s_t, a_t) \left[ q_t(s_t, a_t) - [r_{t+1} + \gamma \max_{a \in A} q_t(s_{t+1}, a)] \right], & \text{if } (s,a) = (s_t, a_t) \\
    q_t(s, a), & \forall (s, a) \neq (s_t, a_t)
\end{cases}
\]

Q-learning解决的是一个贝尔曼最优方程：
$$
q(s, a) = \mathbb{E} \left[ R_{t+1} + \gamma \max_{a'} q(S_{t+1}, a') \, \middle| \, S_t = s, A_t = a \right], \quad \forall s, a.
$$
\subsection*{on-policy 与 off-policy}
behavior policy：用于与环境交互，得到experience的policy

target policy：不断更新，以得到最优的策略

on-policy：两个策略相同，即在与环境交互的过程中改进策略

off-policy：两个策略不用，即使用一个策略与环境交互得到exp，这些exp用来改进另一个策略
\begin{itemize}
    \item MC是一种on-policy的方法
    \item Sarsa是一种on-policy的方法
    \item Q-learning是一种off-policy的方法
\end{itemize}
\subsection{Q-learning的两种形式}
\subsection*{on-policy}
\textbf{1. 更新 Q 值 (Update q-value):}
\[
q_{t+1}(s_t, a_t) = q_t(s_t, a_t) - \alpha_t(s_t, a_t) \left[ q_t(s_t, a_t) - [r_{t+1} + \gamma \max_{a} q_t(s_{t+1}, a)] \right]
\]

\textbf{2. 更新 (行为) 策略 (Update policy):}
\[
\pi_{t+1}(a|s_t) =
\begin{cases}
    1 - \varepsilon + \frac{\varepsilon}{|A(s_t)|} & \text{如果 } a = \underset{a'}{\arg\max} \, q_{t+1}(s_t, a') \\
    \frac{\varepsilon}{|A(s_t)|} & \text{对于其他动作}
\end{cases}
\]
\subsection*{off-policy}
\textbf{1. 更新 Q 值 (Update q-value):}
\[
q_{t+1}(s_t, a_t) = q_t(s_t, a_t) - \alpha_t(s_t, a_t) \left[ q_t(s_t, a_t) - [r_{t+1} + \gamma \max_{a} q_t(s_{t+1}, a)] \right]
\]

\textbf{2. 更新目标策略 (Update target policy):}
\[
\pi_{T, t+1}(a|s_t) =
\begin{cases}
    1 & \text{如果 } a = \underset{a'}{\arg\max} \, q_{t+1}(s_t, a') \\
    0 & \text{对于其他动作}
\end{cases}
\]
\textcolor{blue}{\kaishu*on-policy的方法是$\epsilon$-greedy的，因为它需要用来生成exp，有一定的探索性；而off-policy的策略是greedy的，因为它不需要用来生成exp，但是需要是最优的}

\end{document}